\SectionFrame{3. Platforma predmetu}{Arduino UNO R3 ako doska, ATmega328P ako mikrokontrolér, bare-metal AVR C ako spôsob programovania.}

\begin{frame}{Prečo Arduino UNO R3, keď nechceme Arduino framework}
\begin{itemize}
  \item UNO R3 rieši praktické veci: napájanie, USB, reset, hodinový zdroj, bootloader a konektory.
  \item Tým sa zníži počet chýb spôsobených breadboardom a vodičmi.
  \item Študent sa môže sústrediť na mikrokontrolér, registre a periférie.
  \item Doska je iba nosič ATmega328P, nie dôvod používať \texttt{digitalWrite()}.
\end{itemize}
\begin{alertblock}{Pravidlo predmetu}
Hardvér Arduino UNO R3 áno. Arduino framework nie.
\end{alertblock}
\end{frame}

\begin{frame}{ATmega328P: základné parametre}
\small
\begin{tabular}{p{0.34\textwidth}p{0.53\textwidth}}
\toprule
Parameter & Význam pre predmet \\
\midrule
8-bit AVR RISC & jednoduchá a čitateľná architektúra \\
32 kB Flash & program uložený v mikrokontroléri \\
2 kB SRAM & premenné, zásobník, buffre \\
1 kB EEPROM & trvalé konfiguračné dáta \\
23 GPIO liniek & digitálne vstupy a výstupy \\
3 časovače/čítače & časovanie, PWM, meranie impulzov \\
10-bit ADC & meranie analógových snímačov \\
USART, SPI, TWI & sériová komunikácia s okolím \\
Watchdog & zotavenie z poruchového stavu \\
\bottomrule
\end{tabular}
\end{frame}

\begin{frame}{ATmega328P ako blokový systém}
\begin{columns}[T,onlytextwidth]
\column{0.48\textwidth}
\begin{itemize}
  \item Jadrom je AVR CPU s registrami a ALU.
  \item Periférie sú pripojené cez vnútorné zbernice.
  \item Každá periféria má riadiace, stavové a dátové registre.
  \item Program pracuje s hardvérom práve cez tieto registre.
\end{itemize}
\column{0.47\textwidth}
\includegraphics[width=\linewidth,height=5.2cm,keepaspectratio]{atmega328_block.png}
\end{columns}
\end{frame}

\begin{frame}{Register ako rozhranie medzi programom a hardvérom}
\begin{columns}[T,onlytextwidth]
\column{0.46\textwidth}
\begin{itemize}
  \item Register nie je obyčajná premenná v RAM.
  \item Bit v registri môže povoliť prerušenie, nastaviť režim časovača alebo zmeniť smer pinu.
  \item Datasheet presne definuje význam každého bitu.
\end{itemize}
\column{0.49\textwidth}
\includegraphics[width=\linewidth,height=4.8cm,keepaspectratio]{register_example.png}
\end{columns}
\end{frame}

\begin{frame}[fragile]{Príklad: zápis do registra portu}
\begin{lstlisting}[style=avrstyle]
DDRB  |= (1 << DDB5);    // PB5 ako vystup
PORTB |= (1 << PORTB5);  // log. 1 na PB5
PORTB &= ~(1 << PORTB5); // log. 0 na PB5
\end{lstlisting}
\vspace{0.1cm}
\begin{itemize}
  \item \texttt{DDRB}: smer pinov portu B,
  \item \texttt{PORTB}: výstupná hodnota alebo pull-up,
  \item \texttt{PINB}: čítanie logického stavu pinov.
\end{itemize}
\begin{alertblock}{Most k ďalšej prednáške}
V druhej prednáške už z týchto troch registrov spravíme LED, tlačidlo a jednoduché digitálne rozhranie.
\end{alertblock}
\end{frame}

\begin{frame}[fragile]{Základný program v tomto predmete}
\begin{lstlisting}[style=avrstyle]
#include <avr/io.h>
#include <avr/interrupt.h>

int main(void)
{
    // konfiguracia registrov periferii

    while (1)
    {
        // hlavna slucka programu
    }

    return 0;
}
\end{lstlisting}
\end{frame}

\begin{frame}{Vývojový reťazec}
\begin{center}
\begin{tikzpicture}[x=1cm,y=1cm]
\node[draw=MainBlue,fill=LightBlue,rounded corners=2pt,minimum width=2.3cm,minimum height=0.85cm] (src) at (0,1.6) {\texttt{main.c}};
\node[draw=MainBlue,fill=LightBlue,rounded corners=2pt,minimum width=2.3cm,minimum height=0.85cm] (gcc) at (3.0,1.6) {AVR-GCC};
\node[draw=MainBlue,fill=LightBlue,rounded corners=2pt,minimum width=2.3cm,minimum height=0.85cm] (hex) at (6.0,1.6) {\texttt{.hex}};
\node[draw=MainBlue,fill=MainBlue,text=white,rounded corners=2pt,minimum width=2.3cm,minimum height=0.85cm] (mcu) at (9.0,1.6) {ATmega328P};
\draw[->,MainBlue,line width=1pt] (src) -- (gcc);
\draw[->,MainBlue,line width=1pt] (gcc) -- (hex);
\draw[->,MainBlue,line width=1pt] (hex) -- node[above]{AVRDUDE} (mcu);
\end{tikzpicture}
\end{center}
\BlueBox{Úloha PlatformIO}{PlatformIO je prostredie, ktoré volá prekladač, linker a nahrávací nástroj. Cieľom predmetu však ostáva porozumenie mikrokontroléru.}
\end{frame}

\begin{frame}[fragile]{Minimálny projekt v PlatformIO}
\begin{lstlisting}[style=avrstyle]
; platformio.ini
[env:uno]
platform = atmelavr
board = uno

build_flags =
    -std=gnu11
    -Wall
    -Wextra
    -Wpedantic
\end{lstlisting}
\begin{alertblock}{Kontrola}
V projekte nesmie byť uvedené \texttt{framework = arduino}. Zdrojový súbor má byť \texttt{main.c}, nie \texttt{main.ino}.
\end{alertblock}
\end{frame}

\begin{frame}{Prvé laboratórne cvičenie}
\begin{itemize}
  \item inštalácia a kontrola VS Code, PlatformIO, AVR toolchainu,
  \item vytvorenie projektu pre dosku \texttt{uno},
  \item kompilácia minimálneho programu,
  \item nahratie pripraveného programu do UNO R3,
  \item orientácia na doske: napájanie, USB, reset, digitálne a analógové piny,
  \item vysvetlenie mapovania Arduino pinov na porty ATmega328P.
\end{itemize}
\end{frame}

\begin{frame}{Ako budeme čítať datasheet}
\begin{columns}[T,onlytextwidth]
\column{0.48\textwidth}
\BlueBox{Postup}{
\begin{enumerate}
  \item nájsť kapitolu periférie,
  \item identifikovať registre,
  \item pochopiť význam bitov,
  \item nastaviť bity maskou,
  \item overiť správanie meraním.
\end{enumerate}}
\column{0.48\textwidth}
\BlueBox{Inžinierska disciplína}{
\begin{itemize}
  \item žiadne slepé kopírovanie kódu,
  \item komentovať rozhodujúce bity,
  \item myslieť na elektrické limity,
  \item uvažovať bezpečný stav zariadenia.
\end{itemize}}
\end{columns}
\end{frame}

\begin{frame}{Prepojenie na druhú prednášku}
\begin{itemize}
  \item Dnes sme vysvetlili, prečo mikrokontrolér vznikol a čo obsahuje.
  \item Nabudúce prejdeme ku konkrétnej periférii: digitálne vstupy a výstupy.
  \item Budeme používať registre \texttt{DDRx}, \texttt{PORTx} a \texttt{PINx}.
  \item Prvé praktické úlohy budú LED, tlačidlo, pull-up rezistor, aktívna nula a bitové masky.
\end{itemize}
\begin{alertblock}{Otázka na prípravu}
Prečo je bezpečnejšie nastaviť iba jeden bit registra maskou než zapísať do celého registra novú hodnotu?
\end{alertblock}
\end{frame}
